%%%%%%%%%%%%%%%%%
% CV tailored for SOFTEAM - Ingénieur en Intelligence Artificielle
% Positioning: Ingénieur IA / Data Scientist senior | Python, ML, LLMs, R&D, CIR & acculturation IA
%%%%%%%%%%%%%%%%%

\documentclass[8pt,a4paper,ragged2e,withhyper]{altacv}

\geometry{left=1.25cm,right=1.25cm,top=.5cm,bottom=1.cm,columnsep=1.2cm}

\usepackage{paracol}
\usepackage{fontawesome5}
\usepackage{hyperref}

\iftutex 
  \setmainfont{Roboto Slab}
  \setsansfont{Lato}
  \renewcommand{\familydefault}{\sfdefault}
\else
  \usepackage[rm]{roboto}
  \usepackage[defaultsans]{lato}
  \renewcommand{\familydefault}{\sfdefault}
\fi

\definecolor{SlateGrey}{HTML}{2E2E2E}
\definecolor{LightGrey}{HTML}{666666}
\definecolor{DarkPastelRed}{HTML}{8AC0D9}
\definecolor{PastelRed}{HTML}{00B0DA}
\definecolor{GoldenEarth}{HTML}{B4AD0C}

\colorlet{name}{black}
\colorlet{tagline}{PastelRed}
\colorlet{heading}{DarkPastelRed}
\colorlet{headingrule}{GoldenEarth}
\colorlet{subheading}{PastelRed}
\colorlet{accent}{PastelRed}
\colorlet{emphasis}{SlateGrey}
\colorlet{body}{LightGrey}

\definecolor{violet}{HTML}{5A2582}
\definecolor{rouge}{HTML}{F53B3B}
\definecolor{noir}{HTML}{00232C}
\definecolor{bleu}{HTML}{00B0DA}
\definecolor{rose}{HTML}{F831A0}
\definecolor{jaune}{HTML}{FFEC00}
\definecolor{vert}{HTML}{34AD6C}

\renewcommand{\namefont}{\Huge\rmfamily\bfseries}
\renewcommand{\personalinfofont}{\footnotesize}
\renewcommand{\cvsectionfont}{\LARGE\rmfamily\bfseries}
\renewcommand{\cvsubsectionfont}{\large\bfseries}
\renewcommand{\cvItemMarker}{{\small\textbullet}}
\renewcommand{\cvRatingMarker}{\faCircle}

\input{../_shared/pubs-num.tex}

% auto_tex_manager layout tuning
\emergencystretch=3em
\hfuzz=60pt
\hbadness=9999
\sloppy

\begin{document}

\name{BOILEAU Guillaume}
\tagline{Ingénieur IA / Data Scientist senior | Python, ML, LLMs, R\&D, CIR \& acculturation IA}

\photoL{2.8cm}{../_shared/moi}

\personalinfo{%
  \email{guillaume.boileaupro@gmail.com}
  \phone{+33 6 59 94 85 84}
  \location{Alpes-Maritimes, FRANCE}
  \linkedin{guillaume-boileau-891b81265}
  \researchgate{Guillaume-Boileau-2}
  \orcid{0000-0002-3576-6968}
}

\makecvheader
\vspace{-0.3cm}

\AtBeginEnvironment{itemize}{\small}
\columnratio{0.64}

\begin{paracol}{2}

\cvsection{Experience}

\cvevent{\textbf{Software Engineer -- Développement logiciel, intégration \& IA appliquée}}{VEV Platform Services France}{2025 -- 2026}{Valbonne, France}

\textbf{Contexte :} Environnement industriel logiciel pour une plateforme CPMS liée à des infrastructures de recharge de véhicules électriques, avec services backend, interfaces applicatives, données opérationnelles et équipements terrain.

\textbf{Objectif :} Contribuer au développement, à l’intégration et à la fiabilisation de services logiciels connectés, en combinant développement applicatif, analyse de flux, investigation d’incidents, documentation technique et usages IA pour l’assistance au développement.

\begin{itemize}
  \item \textbf{Développement logiciel :} Contribution à des composants TypeScript, Angular et Node.js intégrés dans une plateforme opérationnelle.
  \item \textbf{Intégration logicielle :} Travail sur les interfaces entre services applicatifs, données opérationnelles et contraintes terrain.
  \item \textbf{Analyse de flux :} Investigation d’échanges FTP, interruptions de flux, incohérences de données et comportements inattendus.
  \item \textbf{Support opérationnel :} Diagnostic technique, analyse d’incidents et contribution aux actions de correction ou de fiabilisation.
  \item \textbf{Documentation technique :} Formalisation de constats, analyses d’anomalies et éléments utiles à la résolution.
  \item \textbf{IA appliquée :} Usage d’outils IA pour relecture de code, debugging, synthèse technique, documentation et structuration de raisonnements.
  \item \textbf{Pratiques outillées :} Utilisation de Git, Linux, Zenhub et workflows collaboratifs pour suivre les sujets techniques.
\end{itemize}

\divider

\cvevent{\textbf{Ingénieur R\&D senior -- IA scientifique, Python \& validation de modèles}}{CNES}{2023 -- 2025}{Nice, France}

\textbf{Contexte :} Activités d’ingénierie et de R\&D pour la mission spatiale LISA ESA/NASA, impliquant simulation, Python scientifique, traitement de données, modélisation statistique, validation technique et coordination internationale.

\textbf{Objectif :} Développer et structurer des workflows Python reproductibles pour intégrer, comparer et valider des modèles hétérogènes, analyser des signaux faibles et produire des résultats traçables pour revue collaborative.

\begin{itemize}
  \item \textbf{Plateforme Python :} Développement de CATpyLISA pour l’intégration, la comparaison, la validation et la revue technique de modèles.
  \textcolor{DarkPastelRed}{\href{https://gitlab.oca.eu/gboileau/lisa_l3_tools}{\bf GitLab: CATpyLISA}}
  \item \textbf{Data science :} Structuration de données, sorties de simulation, métriques de comparaison, contrôles de cohérence et analyses reproductibles.
  \item \textbf{Modélisation statistique :} Analyse de signaux faibles, incertitudes, bruit, comparaison de modèles et validation de résultats.
  \item \textbf{Validation de modèles :} Définition de critères d’acceptabilité, règles de comparaison, vérifications et analyse d’écarts.
  \item \textbf{IA / ML appliqués :} Mobilisation de raisonnements data-driven, métriques, reproductibilité et principes d’évaluation de modèles sur données complexes.
  \item \textbf{Coordination technique :} Interface avec contributeurs scientifiques, logiciels et institutionnels dans un environnement spatial international.
  \item \textbf{Documentation \& acculturation :} Rédaction de supports techniques, partage de connaissances, structuration d’informations et préparation de revues.
  \item \textbf{R\&D / CIR :} Identification de verrous scientifiques et techniques, formalisation des méthodes, hypothèses, résultats et limites.
\end{itemize}

\divider

\cvevent{\textbf{Data Scientist -- Machine learning, bruit \& modélisation statistique}}{Universiteit Antwerpen, Belgium}{2021 -- 2023}{Antwerp, Belgium}

\textbf{Contexte :} Travaux de data science pour instruments de précision et détecteurs de nouvelle génération, combinant données capteurs, bruit environnemental, modélisation statistique, machine learning et validation de performance.

\textbf{Objectif :} Développer des workflows robustes et reproductibles pour identifier, caractériser et réduire des sources affectant la qualité de mesure et la sensibilité de systèmes complexes.

\begin{itemize}
  \item \textbf{Data science :} Développement de workflows Python pour analyse de données, statistiques, performance et validation.
  \item \textbf{Machine learning :} Structuration d’approches de réduction de bruits non stationnaires avec canaux témoins, machine learning et deep learning.
  \item \textbf{Modélisation bayésienne :} Développement de pipelines MCMC pour différentes configurations instrumentales et différents niveaux de bruit.
  \item \textbf{Données capteurs :} Analyse de mesures environnementales, propagation, corrélations spatiales et impacts sur la performance.
  \item \textbf{Évaluation de modèles :} Utilisation de figures de mérite, analyses de sensibilité et contrôles de cohérence pour comparer des configurations.
  \item \textbf{Reporting technique :} Production de résultats traçables, synthèses techniques et recommandations pour parties prenantes internationales.
\end{itemize}

\divider

\cvevent{\textbf{Ingénieur R\&D junior -- Signal, simulation \& calcul scientifique}}{ARTEMIS, Observatoire de la Côte d’Azur}{2018 -- 2021}{Nice, France}

\textbf{Contexte :} Travaux de recherche appliquée liés à l’analyse de signaux faibles, à la simulation numérique, au traitement du signal et à la préparation de missions spatiales.

\textbf{Objectif :} Développer des méthodes d’analyse, automatiser des simulations et produire des résultats scientifiques documentés dans un environnement de recherche international.

\begin{itemize}
  \item \textbf{Signal processing :} Développement de méthodes pour analyser et séparer des signaux faibles et superposés.
  \item \textbf{Simulation numérique :} Développement d’approches de simulation pour environnements complexes et grands jeux de données.
  \item \textbf{Python scientifique :} Mise en place de workflows d’analyse, modélisation, visualisation et validation de résultats.
  \item \textbf{Validation technique :} Vérification de cohérence, comparaison de résultats, études de sensibilité et limites méthodologiques.
  \item \textbf{Rédaction scientifique :} Contribution à des publications, rapports, analyses et supports de présentation.
\end{itemize}

\switchcolumn

\cvsection{Profile}

\textbf{Docteur en physique et ingénieur IA / Data Scientist senior avec une expérience solide en Python, machine learning, modélisation statistique, traitement du signal, validation de modèles, R\&D et coordination technique.}

Positionnement pour ce rôle : contribuer à l’émergence de cas d’usage IA utiles, structurer des démarches IA responsables, accompagner les équipes, participer aux offres IA et valoriser les sujets de recherche via des dossiers CIR.

\begin{itemize}
  \item Plus de 5 ans d’expérience professionnelle en Python, data science, modélisation et validation de workflows complexes.
  \item Solide socle IA/ML : apprentissage supervisé, évaluation de modèles, métriques, validation croisée, incertitudes, deep learning concepts.
  \item Expérience concrète en R\&D, identification de verrous scientifiques/techniques, documentation et synthèse de résultats.
  \item Capacité à vulgariser, former, documenter, fédérer et accompagner des équipes autour de pratiques IA.
  \item Culture logicielle et industrielle : Git, Linux, Docker, CI/CD, API, monitoring/logs, support opérationnel et qualité des livrables.
\end{itemize}

\cvsection{Languages}

\cvskill{English}{4}
\divider
\cvskill{Spanish}{2}
\divider

\medskip

\cvsection{Education}

\cvevent{PhD in Physics}{Université Côte d’Azur}{2018 -- 2021}{Nice, France}

\divider

\cvevent{Selective Graduate Program in Fundamental Physics (Magistère)}{Université Paris-Saclay}{2016 -- 2018}{Orsay, France}

\divider

\cvevent{MSc in Astronomy and Astrophysics}{Observatoire de Paris / Université Paris-Saclay}{2017 -- 2018}{Paris, France}

\divider

\cvevent{BSc in Fundamental Physics}{Université Paris-Saclay}{2015 -- 2016}{Orsay, France}

\cvsection{Professional Training}

\cvevent{Machine Learning in Python with scikit-learn}{FUN MOOC -- Inria}{2026}{France Université Numérique}

\small{
Predictive modelling, supervised learning, model evaluation, cross-validation, metrics, pipelines and practical machine learning workflows with Python and scikit-learn.
}

\divider

\cvevent{Supervised Machine Learning: Regression \& Classification}{DeepLearning.AI -- Andrew Ng}{2020}{Coursera}

\divider

\cvevent{Recherche reproductible : principes méthodologiques pour une science transparente}{FUN MOOC -- Inria}{2025}{France Université Numérique}

\divider

\cvevent{Continuous Professional Training -- Software, Data, AI and Systems}{OpenClassrooms}{2024 -- 2026}{}

\small{
Python, SQL, Machine Learning, AI, Linux, JavaScript, TypeScript, Angular, Node.js, Django, REST APIs, C/C++, web development, algorithms and software engineering fundamentals.
}

\end{paracol}

\cvsection{Projects}

\cvevent{CATpyLISA -- Plateforme Python d’intégration, comparaison et validation}{Mission spatiale LISA ESA/NASA}{2023 -- 2025}{}

Plateforme Python dédiée à l’intégration de modèles, à la structuration de données, à la comparaison de résultats, à la validation technique et à la revue collaborative de workflows scientifiques.

\textbf{Rôle :} développement Python, structuration de données, documentation technique, validation de modèles, analyse d’écarts, synthèse de résultats et coordination avec plusieurs contributeurs.

\textbf{Compétences mobilisées :} Python, NumPy, SciPy, pandas, Matplotlib, validation de modèles, data quality, traitement du signal, statistiques, simulation, GitLab, documentation.

\divider

\cvevent{IA appliquée, agents, chatbots \& acculturation}{LLMs / Agentic AI / MCP / Productivité logicielle}{2024 -- 2026}{}

Exploration et usage de workflows IA pour revue de code, debugging, documentation, synthèse technique, raisonnement incrémental, logique de chatbot et architectures orientées outils.

\textbf{Rôle :} prompting, analyse de cas d’usage, conception de workflows, assistance au développement, documentation, évaluation de pratiques et veille IA.

\textbf{Compétences mobilisées :} LLMs, agents IA, MCP, chatbots, prompting, Python, architecture de workflows, vulgarisation technique, documentation, bonnes pratiques.

\divider

\cvevent{Noise reduction and predictive modelling workflows}{Machine learning / Sensor data / Precision instrumentation}{2021 -- 2023}{}

Développement de workflows numériques et statistiques pour la détection de signaux faibles, l’analyse de données capteurs, la réduction de bruit non stationnaire et la comparaison de configurations système.

\textbf{Rôle :} analyse de données, modélisation statistique, pipelines MCMC, interprétation de données capteurs, approches ML/deep learning pour la réduction de bruit et validation de performance.

\textbf{Compétences mobilisées :} Python, statistiques, MCMC, machine learning, deep learning concepts, traitement du signal, données capteurs, validation, reporting technique.

\cvsection{Compétences clés}

\vspace{-.3cm}

\cvsubsection{IA, data science \& machine learning}

{\small
\cvtag{Intelligence Artificielle}
\cvtag{Machine Learning}
\cvtag{Data Science}
\cvtag{Python}
\cvtag{scikit-learn}
\cvtag{Predictive Analytics}
\cvtag{Modélisation statistique}
\cvtag{Apprentissage supervisé}
\cvtag{Régression}
\cvtag{Classification}
\cvtag{Évaluation de modèles}
\cvtag{Validation croisée}
\cvtag{Métriques}
\cvtag{Feature Engineering}
\cvtag{Deep Learning concepts}
\cvtag{NLP -- ramp-up}
}

\divider\smallskip
\vspace{-.3cm}

\cvsubsection{LLMs, agents \& usages IA}

{\small
\cvtag{LLMs}
\cvtag{Agentic AI}
\cvtag{MCP}
\cvtag{Chatbots}
\cvtag{Prompt Engineering}
\cvtag{Cas d’usage IA}
\cvtag{IA responsable}
\cvtag{Automatisation}
\cvtag{Veille IA}
\cvtag{Outillage IA}
\cvtag{Bonnes pratiques}
\cvtag{Acculturation IA}
\cvtag{Formation}
\cvtag{Communauté IA}
}

\divider\smallskip
\vspace{-.3cm}

\cvsubsection{R\&D, CIR \& offres}

{\small
\cvtag{R\&D}
\cvtag{Crédit Impôt Recherche}
\cvtag{Dossiers CIR -- ramp-up}
\cvtag{Identification de verrous}
\cvtag{Incertitudes scientifiques}
\cvtag{État de l’art}
\cvtag{Méthodologie R\&D}
\cvtag{Rédaction technique}
\cvtag{Synthèse complexe}
\cvtag{Appels d’offres -- support}
\cvtag{AVV -- awareness}
\cvtag{Documentation scientifique}
\cvtag{Vulgarisation}
}

\divider\smallskip
\vspace{-.3cm}

\cvsubsection{Python, software \& industrialisation}

{\small
\cvtag{Python}
\cvtag{NumPy}
\cvtag{SciPy}
\cvtag{pandas}
\cvtag{Matplotlib}
\cvtag{Jupyter}
\cvtag{SQL}
\cvtag{Git}
\cvtag{GitLab}
\cvtag{Linux}
\cvtag{Docker}
\cvtag{CI/CD}
\cvtag{REST API}
\cvtag{TypeScript}
\cvtag{Angular}
\cvtag{Node.js}
\cvtag{OpenSearch}
\cvtag{Sentry}
}

\divider\smallskip
\vspace{-.3cm}

\cvsubsection{Validation, qualité \& systèmes complexes}

{\small
\cvtag{Validation de modèles}
\cvtag{Data Quality}
\cvtag{Reproductibilité}
\cvtag{Traçabilité}
\cvtag{Contrôles de cohérence}
\cvtag{Analyse d’écarts}
\cvtag{Analyse d’anomalies}
\cvtag{Root cause analysis}
\cvtag{Simulation numérique}
\cvtag{Traitement du signal}
\cvtag{Analyse de bruit}
\cvtag{Systèmes complexes}
\cvtag{Support décisionnel}
}

\divider\smallskip
\vspace{-.3cm}

\cvsubsection{Coordination, pédagogie \& outils}

{\small
\cvtag{Coordination transverse}
\cvtag{Accompagnement d’équipes}
\cvtag{Formation interne}
\cvtag{Animation technique}
\cvtag{Documentation}
\cvtag{Confluence}
\cvtag{Jira}
\cvtag{Zenhub}
\cvtag{Power BI}
\cvtag{Agile}
\cvtag{Kanban}
\cvtag{Anglais technique}
\cvtag{Environnement international}
}

\end{document}
